\documentclass[aps,prl,reprint,superscriptaddress,amsmath,amssymb]{revtex4-2}
\usepackage{graphicx,bm,amsthm}
\usepackage[colorlinks=true,citecolor=blue,urlcolor=blue,linkcolor=blue]{hyperref}
\hypersetup{pdftitle={How long does a scrambled quantum memory remember?},pdfauthor={Xi Zhao},pdfsubject={Information-spectrum thresholds for scrambled encodings under local Pauli noise}}
\newtheorem{theorem}{Theorem}
\newtheorem{corollary}{Corollary}
\newcommand{\E}{\mathbb E}
\newcommand{\Tb}{\overline T}
\newcommand{\N}{\mathcal N}
\newcommand{\tr}{\operatorname{tr}}
\newcommand{\prob}{\mathbb P}
\begin{document}
\title{How long does a scrambled quantum memory remember?}
\author{Xi Zhao}\noaffiliation
\date{October 28, 2026}
\begin{abstract}
Scrambling spreads quantum information over high-weight Pauli operators, which
local noise damps exponentially, so scrambled encodings are expected to be
fragile. We show they survive up to a sharp threshold. In an idling memory,
states from any state 2-design remain distinguishable until the Pauli error
entropy reaches one bit per qubit. A random codebook of such states carries
classical information at exactly the rate $1-H$, and beyond the threshold a
single noise layer forces exponential error-mitigation cost. Purity bounds
overestimate the memory lifetime almost twofold.
\end{abstract}
\maketitle

\emph{Introduction.}---Quantum information is rarely stored in simple states.
A quantum memory, a register waiting between gates, or the output of a generic
circuit holds a delocalized, scrambled state, modeled by a state 2-design: an
ensemble, such as random stabilizer states, whose first two moments match
Haar-random states~\cite{KuengGross}. While it waits, every qubit decoheres.
How long does the stored information last?

The standard picture suggests it does not last long. A scrambled state places
almost all of its Pauli weight on operators that act on a constant fraction of
the qubits, and local depolarizing noise of strength $r$ multiplies a weight-$w$
operator by $(1-r)^w$. At $r=0.2$ and $n=100$ qubits, the surviving Pauli
weight of a 2-design state is about $10^{-14}$. The same mechanism underlies the
exponential sample cost of quantum error mitigation (QEM). Takagi, Tajima, and
Gu~\cite{Takagi} and Tsubouchi, Sagawa, and Yoshioka~\cite{Tsubouchi} proved
that this cost grows exponentially with circuit depth for any noise strength;
related analyses appear in Refs.~\cite{TakagiPRR,TakagiNPJ}. Quek
\emph{et al.}~\cite{Quek} showed that random circuits reach this regime at very
shallow depth, using purity to control the decay.

Pauli weight, however, is not information. We show that a scrambled memory
keeps its information until a sharp threshold, fixed by the entropy of the
single-qubit errors, and we identify three consequences. First, a memory with
coherence times $T_1=T_2$ retains information for $0.291\,T_1$, whereas purity
certifies loss only after $0.549\,T_1$. Second, a codebook of scrambled states
transmits classical information at exactly the hashing rate $1-H$. Third, beyond
the threshold, a single layer of noise already forces QEM to use exponentially
many samples. For biased noise and pure dephasing, the gap to purity-based
estimates widens qualitatively.

All three follow from one quantity. Two pure states $\psi_1,\psi_2$ are drawn
independently from a state 2-design and stored under the same Pauli channel
$\N$, with no further randomization. With equal priors, the optimal probability
of identifying the stored state is $(1+\Tb)/2$, where
\begin{equation}
\Tb=\E_{\psi_1,\psi_2}\tfrac12\|\N(\psi_1)-\N(\psi_2)\|_1 .
\label{eq:observable}
\end{equation}
The absence of randomization after the noise is essential. Averaging
$U^\dagger\N(U\cdot U^\dagger)U$ over a unitary 2-design turns any Pauli channel
into global depolarization, with parameter
$\theta=(d^2p_I-1)/(d^2-1)$~\cite{Dankert}, which decays exponentially in $n$ for
every $r>0$. At $n=8$ and $r=0.2$, $\Tb=0.736$ without the twirl, against about
$0.272$ with it [Fig.~\ref{fig:bounds}(a)].

Ibarrondo and Stilck Fran\c{c}a~\cite{Ibarrondo} and Liu \emph{et al.}~\cite{Liu}
bounded $\Tb$ for local depolarizing noise, finding near-perfect
distinguishability below $r\approx0.25$ and decay above $r\approx0.42$, with the
interval between them open. Sommers \emph{et al.}~\cite{Sommers} found the hashing
transition, with an $n^{-1/2}$ window, in the coherent information of random
codes. The bound below closes the interval and holds for every system size.

\emph{Information-spectrum bracket.}---Let $d=2^n$,
$\mathcal P_n=\{I,X,Y,Z\}^{\otimes n}$, and
\begin{equation}
\N_p(\rho)=\sum_{P\in\mathcal P_n}p_PP\rho P,\qquad s(P)=-\log_2p_P,
\label{eq:channel}
\end{equation}
with $p$ an arbitrary distribution, so errors may be correlated, and with the
surprisal $s$ distributed as $P\sim p$. Logarithms are base two. Each
$\psi_i=|\psi_i\rangle\langle\psi_i|$ is an $n$-qubit pure state, generally
entangled, drawn independently from an exact state 2-design on the full space,
\begin{equation}
\E\psi_i=I/d,\qquad\E\psi_i^{\otimes2}=(I+F)/[d(d+1)],
\label{eq:design}
\end{equation}
where $F$ swaps the two copies; the two ensembles may differ.

\begin{theorem}
For every Pauli distribution $p$ and every real $a$,
\begin{equation}
\prob[s<n-a]-2^{-a}\le\Tb\le\prob[s<n+a]+2^{-a/2}/\sqrt2 .
\label{eq:spectrum}
\end{equation}
More generally, for any set $\mathcal G\subseteq\mathcal P_n$ and any partition
$\{\mathcal C_j\}$ of $\mathcal P_n$, with $m_j=\sum_{\mathcal C_j}p_P$ and
$v_j=\sum_{\mathcal C_j}p_P^2$,
\begin{align}
\Tb&\ge p(\mathcal G)-|\mathcal G|/d,\label{eq:lower}\\
\Tb&\le\sum_j\min\left\{m_j,\sqrt{\frac{d^2v_j-m_j^2}{2(d+1)}}\right\}.
\label{eq:upper}
\end{align}
\end{theorem}

\begin{proof}
Fix $\psi_1$ and let $\Pi$ project onto the span of
$\{P|\psi_1\rangle:P\in\mathcal G\}$. For each realization,
$\frac12\|\N_p(\psi_1)-\N_p(\psi_2)\|_1\ge\tr\Pi[\N_p(\psi_1)-\N_p(\psi_2)]$.
The first term is at least $p(\mathcal G)$. Averaged over the independent
$\psi_2$, the second is $\tr\Pi/d\le|\mathcal G|/d$, using only
$\E\psi_2=I/d$. Averaging over $\psi_1$ gives~\eqref{eq:lower}, the
typical-image argument of Ref.~\cite{Ibarrondo} applied to an arbitrary set of
Pauli errors. For the upper bound let $\Delta=\psi_1-\psi_2$ and
$\N_j(\rho)=\sum_{\mathcal C_j}p_PP\rho P$. Pauli orthogonality
and~\eqref{eq:design} give
\begin{align}
\E\tr\N_j(\psi_i)^2&=\frac{dv_j+m_j^2}{d+1},\qquad
\E\tr\N_j(\psi_1)\N_j(\psi_2)=\frac{m_j^2}{d},\nonumber\\
\E\tr\N_j(\Delta)^2&=\frac{2d}{d+1}v_j-\frac{2m_j^2}{d(d+1)}.
\label{eq:moment}
\end{align}
Jensen's inequality and $\|A\|_1\le\sqrt d\|A\|_2$ bound
$\E\frac12\|\N_j(\Delta)\|_1$ by the square root in~\eqref{eq:upper}, and
$\|\Delta\|_1\le2$ bounds it by $m_j$. For~\eqref{eq:spectrum}, take
$\mathcal G=\{s<n-a\}$, which has fewer than $d2^{-a}$ elements, and split the
partition at $s=n+a$, where the tail has $v\le2^{-n-a}$.
\end{proof}

The theorem has a simple reading. The error images $P|\psi_1\rangle$ of a
scrambled state are nearly orthogonal, and a $d$-dimensional space holds at
most $d$ of them. Errors with $p_P>1/d$, i.e., $s(P)<n$, number fewer than $d$,
so their images fit into the space, and the projector onto these images
separates $\psi_1$ from an independent $\psi_2$. The remaining errors are too
numerous to fit, and their contributions average out. Distinguishability is
therefore pinned between the probabilities of a likely error at surprisal
$n\mp a$, up to exponentially small corrections. We checked this argument in
Lean~4 with Mathlib~\cite{Lean,Mathlib}: from Eq.~\eqref{eq:lower} and the
two-class case of Eq.~\eqref{eq:upper} as hypotheses, the formal proof derives
Eq.~\eqref{eq:spectrum} and Corollary~1 below, and it proves the counting and
moment algebra in full~\cite{DataRepo}.

The trivial partition, $\mathcal C=\mathcal P_n$, reduces~\eqref{eq:upper} to the
purity of the noisy output, $\Tb\le[d\sum_Pp_P^2/2]^{1/2}$, the quantity that
controls the bounds of Ref.~\cite{Quek}. Purity is loose because it places
likely and unlikely errors in the same second moment. Separating them at
surprisal $n$ removes this looseness.

\emph{Sharp threshold.}---For i.i.d. noise, $p=\bm p^{\otimes n}$, with
single-qubit distribution $\bm p=(p_0,p_x,p_y,p_z)$, entropy
$H=\E[-\log_2p_i]$, and variance $\sigma^2=\operatorname{Var}[-\log_2p_i]$. The
surprisal is then a sum of $n$ independent terms with mean $nH$. Taking
$a=\epsilon n$ in~\eqref{eq:spectrum} and applying a large-deviation
bound~\cite{CoverThomas} gives
\begin{equation}
\Tb=\begin{cases}1-e^{-\Omega(n)},&H<1,\\ e^{-\Omega(n)},&H>1.\end{cases}
\label{eq:threshold}
\end{equation}
The threshold $H=1$ is the hashing point, at which the coherent information
$1-H$ of the channel with maximally mixed input vanishes~\cite{Bennett,Devetak}.
For depolarizing noise $\N_r(\rho)=(1-r)\rho+rI/2$, it lies at
$r_c=0.252386\ldots$, whereas purity certifies loss only for
$r>1-1/\sqrt3=0.423$. Above threshold the decay rate is at least the
universal$_2$ privacy-amplification exponent~\cite{Hayashi,Dupuis}; near
$r_c$ both upper and lower rates equal $2.64\,(r-r_c)^2$ to leading
order~\cite{DataRepo}.

The bracket also resolves the shape of the transition. Let $\bm p_n$ depend on
$n$, with i.i.d. errors at each $n$, and set $x_n=\sqrt n(1-H_n)/\sigma_n$.

\begin{corollary}
If $\bm p_n\to\bm p_*$ with $H(\bm p_*)=1$ and $\sigma(\bm p_*)>0$, then for
every pair of independent exact 2-design ensembles,
$\Tb_n=\Phi(x_n)+O(\log n/\sqrt n)$, where $\Phi$ is the standard normal
distribution function.
\end{corollary}

\begin{proof}
The Berry--Esseen theorem~\cite{BerryEsseen} gives
$\prob[s<n\mp a]=\Phi(x_n\mp a/\sqrt n\sigma_n)+O(n^{-1/2})$, uniformly because
the third moment of the surprisal is bounded on a four-letter alphabet and
$\sigma_n$ is bounded away from zero. Insert into~\eqref{eq:spectrum} with
$a=2\log_2n$.
\end{proof}

The variance condition is necessary. Complete dephasing has $H=1$ but
$\sigma=0$. The union of $d+1$ mutually unbiased bases is an exact
2-design~\cite{MUB}, and with the computational basis among them,
$\Tb_{\rm MUB}=(d-1)(2d+1)/[d(d+1)^2]\to0$, whereas
$\Tb_{\rm Haar}=(d-1)/(2d-1)\to1/2$. Two ensembles with the same two moments
reach different critical limits, sharpening an observation of
Ref.~\cite{Ibarrondo}.

\begin{figure*}[t]
\includegraphics[width=\textwidth]{figures/bounds.pdf}
\caption{Distinguishability after one layer of local depolarizing noise.
(a) At $n=8$, the shaded band lies between the projector lower bound (solid)
and the two-class upper bound (dashed) of Eqs.~\eqref{eq:lower}
and~\eqref{eq:upper}, optimized over error-weight cutoffs; the thin vertical
line marks $r_c$. Points are averages over 64 independent Haar-state pairs from
full density matrices, with one-standard-error bars; the same pairs are used
at every $r$. The dash-dotted curve is the purity bound, which drops below one
only for $r\gtrsim0.37$ at this size and certifies decay only above $0.423$
asymptotically. The dotted gray curve is the depolarizing parameter $\theta$
reached if a further 2-design unitary twirled the noise; it equals the
twirled distance up to a factor $1-O(1/d)$. (b) The bracket~\eqref{eq:spectrum}
with $a=2\log_2n$, evaluated from binomial distributions at fixed
$x=\sqrt n(1-H)/\sigma$ for $n=1024$ and $65536$, approaches $\Phi(x)$ (black).}
\label{fig:bounds}
\end{figure*}

\emph{Memory lifetime.}---An idling qubit with relaxation time $T_1$ and
dephasing time $T_2\le2T_1$ is described, in the Pauli-twirl approximation
standard in threshold analyses~\cite{Ghosh}, by
$p_x=p_y=(1-e^{-t/T_1})/4$ and $p_z=(1-e^{-t/T_2})/2-p_x$. The scrambled memory
retains its information until $H(t_c)=1$. For $T_2=T_1$ the noise is
depolarizing, $r(t)=1-e^{-t/T_1}$, so
\begin{equation}
t_c=-\ln(1-r_c)\,T_1=0.291\,T_1,
\label{eq:lifetime}
\end{equation}
whereas purity certifies loss only after $(\ln3/2)T_1=0.549\,T_1$. Across
$0.2\le T_2/T_1\le2$, the purity estimate exceeds $t_c$ by a factor $1.77$ to
$1.89$ (Table~\ref{tab:lifetime}). For pure dephasing ($T_1\to\infty$),
$p_z<1/2$ and hence $H<1$ at every finite $t$: the memory never loses its
information, although for each fixed $t$ the convergence rate of
$\Tb_n\to1$ decreases as $t$ grows. Purity certifies nothing at all in this
case, because $2(p_0^2+p_z^2)\ge1$. These statements are theorems for the
twirled channel; for the untwirled amplitude-damping channel, which is not
unital, they are approximations.

\begin{table}[t]
\caption{Retention time of a scrambled memory, in units of $T_1$, from the
hashing condition $H=1$ and from the purity bound, for Pauli-twirled idling.}
\label{tab:lifetime}
\begin{ruledtabular}
\begin{tabular}{cccc}
$T_2/T_1$ & $t_c$ (hashing) & purity & ratio\\
\hline
2.0 & 0.497 & 0.881 & 1.77\\
1.0 & 0.291 & 0.549 & 1.89\\
0.5 & 0.187 & 0.347 & 1.85\\
0.2 & 0.105 & 0.186 & 1.77
\end{tabular}
\end{ruledtabular}
\end{table}

Biased noise shows a sharper contrast. Let the total error probability be $p$
and the bias $\eta=p_z/(p_x+p_y)$. The hashing threshold rises with bias, from
$p_c=0.189$ at $\eta=1/2$ (depolarizing) to $0.278$ at $\eta=10$ and $0.390$ at
$\eta=100$. Purity, by contrast, certifies loss only on a bounded window of
$p$: for $\eta=10$ only $0.387<p<0.705$, and for $\eta=100$ only
$0.455<p<0.555$. Strong biased noise concentrates on few Pauli errors, which
raises the purity, so purity bounds fail precisely where biased-noise
hardware~\cite{Puri} operates.

\emph{Information capacity of scrambled codebooks.}---A memory that stores one
of $M$ messages in states $\psi_1,\ldots,\psi_M$, drawn independently from a
state 2-design, is a classical code. Its rate is set by the same spectrum.

\begin{theorem}
With the square-root measurement~\cite{HayashiNagaoka}, the average error
probability satisfies
$\E P_{\rm err}\le2\prob[s\ge n-a]+4(M-1)2^{-a}$; this needs only a 1-design.
For every measurement and every $\gamma,a>0$, the average success probability
satisfies
\begin{equation}
\E P_{\rm succ}\le\frac\gamma M+\frac1{4\gamma}+\frac{2^{-a}}4
+\prob[s<n-\log_2\gamma+a],
\label{eq:converse}
\end{equation}
which uses the second moments of the design.
\end{theorem}

\begin{proof}
For the first bound, take $\Pi_i$ as above with $\mathcal G=\{s<n-a\}$ and
apply the Hayashi--Nagaoka inequality
$I-E_i\le2(I-\Pi_i)+4\sum_{j\ne i}\Pi_j$~\cite{HayashiNagaoka}. Then
$\tr\N(\psi_i)\Pi_i\ge p(\mathcal G)$, and independence gives
$\E\tr\N(\psi_i)\Pi_j=\tr\Pi_j/d<2^{-a}$ for $j\ne i$. For the second, any POVM
satisfies $\tr E_i\N(\psi_i)\le\tr(\N(\psi_i)-cI)_++c\tr E_i$, so
$P_{\rm succ}\le cd/M+M^{-1}\sum_i\tr(\N(\psi_i)-cI)_+$. Split $\N(\psi_i)$ at
$p_P=c'$. The head contributes at most its mass. For the tail $B$, use
$(\lambda-c)_+\le\lambda^2/4c$ and
$\E\tr B^2=(dv_t+m_t^2)/(d+1)\le c'+1/d$ from~\eqref{eq:design}, then set
$\gamma=cd$ and $c'=c2^{-a}$.
\end{proof}

For i.i.d. noise and $M=2^{nR}$, the first bound vanishes for $R<1-H$, and
Eq.~\eqref{eq:converse} with $\gamma=2^{n(R-\epsilon)}$ and $a=\epsilon n$
forces $\E P_{\rm succ}\to0$ for $R>\max(1-H,0)$. Scrambled codebooks therefore
carry classical information at exactly the hashing rate $1-H$, with a strong
converse. The 2-design property is needed for the converse: computational-basis
codewords, a 1-design, are blind to $Z$ errors and can exceed $1-H$. Product
inputs with joint decoding achieve the capacity $1-h_2(1-r/2)$ of the
depolarizing channel~\cite{King}, which stays positive for all $r<1$. Throughout
$r_c<r<1$, scrambling the encoding therefore costs the entire classical
capacity; at $r_c$ that capacity is $0.453$ bits per qubit.

\emph{Error mitigation after one noisy layer.}---Suppose a QEM protocol receives
$N$ copies of $\N(\rho)$ and must estimate $\tr(A\rho)$ to accuracy $\delta$
with probability at least $1-\epsilon$, for all pure inputs and all observables
with $\|A\|\le1$. Two inputs whose expectation values differ by more than
$2\delta$ must then be distinguished from $N$ copies, so
$1-2\epsilon\le N\cdot\frac12\|\N(\psi_1)-\N(\psi_2)\|_1$~\cite{Takagi}. By
Markov's inequality, a random 2-design pair has distance at most $2\Tb$ with
probability at least $1/2$, and overlap $|\langle\psi_1|\psi_2\rangle|^2\le4/d$
with probability at least $3/4$. Taking $A=\psi_1$ then gives
\begin{equation}
N\ge\frac{1-2\epsilon}{2\Tb_n}
\label{eq:qem}
\end{equation}
whenever $2\delta<1-4/d$. Beyond the hashing point, $\Tb_n=e^{-\Omega(n)}$, so a
single noisy layer already forces exponentially many samples. For one layer,
Theorem~3 of Ref.~\cite{Takagi} gives a bound that decreases with $n$, and the
purity route of Ref.~\cite{Quek} applies only for $r>0.423$. Below threshold,
Eq.~\eqref{eq:qem} gives no exponential bound. The bound concerns protocols
that only receive copies of the noisy state and that work for global
observables.

\emph{Discussion.}---A scrambled encoding is not fragile under local Pauli noise.
Its information survives until the error entropy reaches one bit per qubit,
and this one threshold controls three quantities at once: how long a scrambled
memory remembers, how much classical information a scrambled codebook carries,
and when a single noisy layer makes error mitigation exponentially expensive.
In each case, purity-based reasoning misplaces the boundary, by nearly a factor
of two in memory lifetime and by an unbounded amount for biased noise and pure
dephasing. The reason is the same throughout: purity weighs the few likely
errors like the many unlikely ones, whereas the optimal measurement resolves
the likely ones.

The bracket rests only on the error distribution and the dimension, so it holds
for every 2-design and for correlated noise. For correlated errors, the full
information spectrum replaces the entropy: global depolarization has entropy
rate $2(1-\theta)>1$ for $\theta<1/2$, yet $\Tb\to\theta$, which the bracket
reproduces. The main open problem is depth. In deeper circuits, later
scrambling layers partly twirl earlier noise toward global depolarization.
Tracking the information spectrum layer by layer would decide whether a finite
threshold survives at depth. Determining Haar decay rates away from threshold,
where random stabilizer states, a 3-design, decay faster than the achievable
exponent~\cite{DataRepo}, is a second open question.

\emph{Acknowledgments.}---Generative AI was used substantively in this work.
Claude Opus~5.5 (Anthropic), accessed through Claude Code in October 2026, with
review passes by other Claude model instances, was directed by the author to
select the problem, draft and revise the derivations, survey the literature,
write the analysis code and the Lean formalization, generate the figures and
table, and draft the text. A review with an OpenAI GPT model [version to be
specified by the author] identified the omitted Ref.~\cite{Ibarrondo}. The
author is responsible for all content. Proofs were rederived independently,
Theorem~1 and Corollary~1 were machine-checked in Lean as described in the
text, the exact counterexamples were checked by explicit enumeration for
$n\le2$, the codebook bounds were checked against semidefinite-programming
optima for $n=2$, numerical claims were checked against independent
implementations recorded with the data, and cited statements were checked
against the primary sources.

\emph{Data availability.}---The data that support the findings of this article,
the code that generates every figure, table, and numerical check, and the Lean
formalization accompany this submission and will be deposited in a public
repository~\cite{DataRepo}.

\begin{thebibliography}{99}
\bibitem{KuengGross} R. Kueng and D. Gross, Qubit stabilizer states are
complex projective 3-designs,
\href{https://arxiv.org/abs/1510.02767}{arXiv:1510.02767} (2015).
\bibitem{Takagi} R. Takagi, H. Tajima, and M. Gu, Universal sampling lower
bounds for quantum error mitigation, Phys. Rev. Lett. \textbf{131}, 210602 (2023).
\bibitem{Tsubouchi} K. Tsubouchi, T. Sagawa, and N. Yoshioka, Universal cost
bound of quantum error mitigation based on quantum estimation theory,
Phys. Rev. Lett. \textbf{131}, 210601 (2023).
\bibitem{TakagiPRR} R. Takagi, Optimal resource cost for error mitigation,
Phys. Rev. Res. \textbf{3}, 033178 (2021).
\bibitem{TakagiNPJ} R. Takagi, S. Endo, S. Minagawa, and M. Gu, Fundamental
limits of quantum error mitigation, npj Quantum Inf. \textbf{8}, 114 (2022).
\bibitem{Quek} Y. Quek, D. Stilck Fran\c{c}a, S. Khatri, J. J. Meyer, and
J. Eisert, Exponentially tighter bounds on limitations of quantum error
mitigation, Nat. Phys. \textbf{20}, 1648 (2024).
\bibitem{Dankert} C. Dankert, R. Cleve, J. Emerson, and E. Livine, Exact and
approximate unitary 2-designs and their application to fidelity estimation,
Phys. Rev. A \textbf{80}, 012304 (2009).
\bibitem{Ibarrondo} R. Ibarrondo and D. Stilck Fran\c{c}a, Average contraction
coefficients of quantum channels,
\href{https://arxiv.org/abs/2508.08214}{arXiv:2508.08214} (2025).
\bibitem{Liu} G. Liu, C. Qin, Z. Xu, X. Ma, and Z.-W. Liu,
Can scrambling protect quantum state distinguishability under noise?,
\href{https://arxiv.org/abs/2606.02122}{arXiv:2606.02122} (2026).
\bibitem{Sommers} G. M. Sommers, J. A. Jacoby, Z. Weinstein, D. A. Huse,
and S. Gopalakrishnan, Spectral properties and coding transitions of
Haar-random quantum codes, PRX Quantum \textbf{7}, 020328 (2026).
\bibitem{Lean} L. de Moura and S. Ullrich, The Lean~4 theorem prover and
programming language, in \emph{Automated Deduction -- CADE 28}, Lecture Notes in
Computer Science Vol.~12699 (Springer, Cham, 2021), pp.~625--635.
\bibitem{Mathlib} The mathlib Community, The Lean mathematical library, in
\emph{Proceedings of the 9th ACM SIGPLAN International Conference on Certified
Programs and Proofs} (ACM, New York, 2020), pp.~367--381.
\bibitem{DataRepo} X. Zhao, Data, code, and Lean formalization for ``How long
does a scrambled quantum memory remember?'' [repository and DOI to be assigned
on deposit].
\bibitem{CoverThomas} T. M. Cover and J. A. Thomas, \emph{Elements of
Information Theory}, 2nd ed. (Wiley, Hoboken, 2006), Ch.~11.
\bibitem{Bennett} C. H. Bennett, D. P. DiVincenzo, J. A. Smolin, and
W. K. Wootters, Mixed-state entanglement and quantum error correction,
Phys. Rev. A \textbf{54}, 3824 (1996).
\bibitem{Devetak} I. Devetak, The private classical capacity and quantum
capacity of a quantum channel, IEEE Trans. Inf. Theory \textbf{51}, 44 (2005).
\bibitem{Hayashi} M. Hayashi, Tight exponential analysis of universally
composable privacy amplification and its applications,
IEEE Trans. Inf. Theory \textbf{59}, 7728 (2013).
\bibitem{Dupuis} F. Dupuis, Privacy amplification and decoupling without
smoothing, IEEE Trans. Inf. Theory \textbf{69}, 7784 (2023).
\bibitem{BerryEsseen} W. Feller, \emph{An Introduction to Probability Theory and
Its Applications}, Vol.~II, 2nd ed. (Wiley, New York, 1971), Sec.~XVI.5.
\bibitem{MUB} A. Klappenecker and M. R\"otteler, Mutually unbiased bases
are complex projective 2-designs, in \emph{Proceedings of the 2005 IEEE
International Symposium on Information Theory} (IEEE, 2005), pp.~1740--1744.
\bibitem{Ghosh} J. Ghosh, A. G. Fowler, and M. R. Geller, Surface code with
decoherence: An analysis of three superconducting architectures,
Phys. Rev. A \textbf{86}, 062318 (2012).
\bibitem{Puri} S. Puri, L. St-Jean, J. A. Gross, A. Grimm, N. E. Frattini,
P. S. Iyer, A. Krishna, S. Touzard, L. Jiang, A. Blais, S. T. Flammia, and
S. M. Girvin, Bias-preserving gates with stabilized cat qubits, Sci. Adv.
\textbf{6}, eaay5901 (2020).
\bibitem{HayashiNagaoka} M. Hayashi and H. Nagaoka, General formulas for
capacity of classical-quantum channels, IEEE Trans. Inf. Theory \textbf{49},
1753 (2003).
\bibitem{King} C. King, The capacity of the quantum depolarizing channel,
IEEE Trans. Inf. Theory \textbf{49}, 221 (2003).
\end{thebibliography}
\end{document}
